\begin{lemma}[The closed-core elementary escape]
Let \(\mathcal E:\noderef{ExtendedTashkinovTree}(G,k,e_0,\rho,T,\varphi)\)
have terminal kind PE in a finite loopless multigraph, and let \(T_n\) be
its terminal ladder tree. Let \(C\) be any ordered tree-sequence reached
from \(T_n\) by \noderef{TashkinovAugmentationSequence} under \(\varphi\).
If \(V(C)\) is elementary and closed under \(\varphi\), then \(V(T)\)
is elementary under \(\varphi\).
\end{lemma}
\begin{proof}
Put \(F=\mathcal E.\mathit{series.data.tree}(\operatorname{last}(n+1))\),
where \(n=\mathcal E.\mathit{rounds}\).
By \noderef{P2634_T1_PETerminalFacts}, \(F\) is a closure of \(T_n\)
under \(\varphi\). The assumed augmentation route and closedness of \(C\)
make \(C\) another such closure, by \noderef{TashkinovClosure}.
Thus \noderef{TashkinovClosureVertexSetUnique} gives \(V(C)=V(F)\).
The below-final field of \noderef{ExtendedTashkinovTree} says that
\(T\preccurlyeq F\) as exact ordered tree-sequences.
By \noderef{TreeSequenceInitialSegmentVertexSetSubset},
\(V(T)\subseteq V(F)=V(C)\). For two distinct vertices of \(V(T)\),
this containment places both in \(V(C)\); the assumed elementarity of
\(V(C)\) gives disjoint missing-color sets. This is precisely
\noderef{ColoredVertexSetElementary} for \(V(T)\).
\end{proof}
